\scanpage{4}
Bien entendu, (3) et (4) sont compatibles avec l'action de $\unsure{\mathfrak S}_I$ (qui opère sur $\omega(I)$ par la signature).

Je définis $S_0\mathcal T_{0,I}$ comme le \unsure{ss-groupoïde} du groupoïde fondamental de $P\mathbb U_{0,I}$ induit \deleted{\unsure{par}} sur $\omega(I)$ comme ens. d'objets. Je voudrais en donner une expression purement algébrique. Pour ceci, soit $\widetilde{\mathbb U}\simeq\mathbb R$ le \inserted{groupe} revêtement universel standard de $\mathbb U$, et considérons la suite exacte
\[
(5_0)\qquad 0-\mathbb Z\longrightarrow
\underset{t\;\mapsto\;\exp 2i\pi t}{\mathbb R\longrightarrow\mathbb U}
\longrightarrow0
\]
d'où
\[
(5_I)\qquad0\longrightarrow\mathbb Z^I\longrightarrow\mathbb R^I\longrightarrow\mathbb U^I\longrightarrow0,
\]
\deleted{\unsure{P}} Notons que $P\mathbb U_{0,I}$ est un torseur sous le groupe topologique $\mathbb U^I$, \deleted{ce qui suffit à}\deleted{\unsure{dét}}\inserted{structure algébrique} qui détermine sa topologie. Je vais donner une construction purement algébrique du groupoïde fondamental de cet espace, en utilisant seulement $(5_I)$ et sa structure de torseur. De façon générale, considérons une structure d'extension de groupes topologiques comme
\[
(6)\qquad0\longrightarrow\pi\longrightarrow\widetilde T\longrightarrow T\longrightarrow0
\]
avec \inserted{$\pi$ discret,} $\pi\subset\widetilde T^0$ et $\widetilde T^0$ simplement connexe
\mspage{3}
